\documentclass[conference]{../../../../Conference-LaTeX-template_10-17-19/IEEEtran}
\usepackage{cite}
\usepackage{amsmath,amssymb,amsfonts}
\usepackage{array}
\usepackage{booktabs}
\usepackage{tabularx}
\usepackage{url}
\newcolumntype{Y}{>{\raggedright\arraybackslash}X}
\renewcommand{\arraystretch}{1.12}

\title{What Is the Ultimate Statistical Invariance of Turbulence?\\A Conditional Universality Research Proposal}
\author{\IEEEauthorblockN{Anonymous Research Proposal}}

\begin{document}
\maketitle

\begin{abstract}
Turbulence is often described as universal because energy cascades, scaling ranges, and normalized statistics recur across apparently unrelated flows. Yet no single measured statistic, exponent, or probability-density collapse is known to be invariant across every turbulent regime. The phrase ``ultimate statistical invariance'' instead conceals several distinct notions: exact symmetries of the governing equations, exact conditional balance relations, asymptotic similarity hypotheses, and empirical cross-flow regularities. This paper develops the \emph{Conditional Universality and Regime Audit for Turbulence} (CURA-Turb), a source-bounded, design-only protocol for separating these claims before they are compared. The protocol represents each assertion as a claim dossier containing a statistic, regime card, scale interval, estimator, uncertainty route, and evidence tier. It uses analytic symmetry references; forced and decaying homogeneous-isotropic cards; wall, rotating/stratified, and sheared/inhomogeneous cards; and an experimental-representation card. The exact Kolmogorov third-order relation is treated as a powerful conditional benchmark, whereas higher-order scaling and distributional collapse are evaluated as potential conditional universality or intermittency-modified behavior. CURA-Turb records exact-law residual intervals, structure-function exponent intervals, normalized PDF distances, anisotropy and inhomogeneity indicators, time-irreversibility indicators, and status stability. It authorizes only conditional labels: exact-relation eligible, conditionally universal, intermittency modified, regime dependent, representation invalid, inconclusive, or model invalid. The study design contains no new direct numerical simulation, laboratory measurement, or engineering-flow run. Its contribution is a falsifiable reporting language that treats universality as a testable validity domain rather than a one-number property of turbulence.
\end{abstract}

\begin{IEEEkeywords}
turbulence, universality, statistical invariance, Kolmogorov scaling, four-fifths law, intermittency, structure functions, Reynolds number, uncertainty quantification
\end{IEEEkeywords}

\section{Introduction}
Turbulence appears in atmospheric flows, oceans, combustion, aircraft wakes, process equipment, energy systems, biological transport, and numerical models. Its visual diversity creates a scientific temptation: perhaps sufficiently small scales forget the geometry and forcing that created them, leaving a single statistical law. That aspiration has motivated a deep collaboration among mathematics, physics, computer science, and engineering. It has also created a recurring ambiguity. A symmetry of the Navier--Stokes equations, an exact inertial-range balance, a fitted scaling exponent, and a collapsed normalized probability density are not interchangeable forms of evidence.

The most defensible answer to the topic is therefore direct. There is no established single ``ultimate statistical invariant'' shared without qualification by all turbulent flows. There are exact statements within defined assumptions, influential similarity hypotheses, and repeatable statistical regularities within compatible flow classes. Their usefulness is not reduced by being conditional; the conditions determine where a claim can be used for theory, model validation, or engineering transfer.

The publisher-verified review by Sreenivasan and Antonia frames small-scale turbulence through classical universality, intermittency, refined similarity, anomalous scaling exponents, and homogeneous-turbulence simulations \cite{sreenivasan_antonia_1997}. That scope already signals a key point: universality is a research program with competing diagnostics, not a single settled scalar. Kolmogorov's 1941 phenomenology remains a natural baseline \cite{kolmogorov_1941}, while the 1962 refinement and later multifractal/cascade accounts make intermittency central rather than peripheral \cite{kolmogorov_1962,frisch_1995}. The present proposal turns those distinctions into a predeclared computational design.

This paper simulates the repository's four-agent workflow:
\begin{center}
Survey $\rightarrow$ Idea $\rightarrow$ ExperimentDesign $\rightarrow$ Author.
\end{center}
The Survey Agent freezes sources, claim boundaries, and research gaps. The Idea Agent selects CURA-Turb rather than an exponent-only score. The ExperimentDesign Agent defines compatible regime cards and a validation ladder. The Author Agent writes only from these upstream artifacts and keeps the document in the status \texttt{PROPOSAL\_NO\_OBSERVED\_RESULTS}. The report supplies a concrete path to stronger turbulence claims without presenting an unrun calculation as evidence.

\section{What Statistical Invariance Can Mean}
\subsection{Equation symmetries are structural, not automatically statistical}
For incompressible flow, a standard velocity-pressure form of the governing equations is
\begin{equation}
\partial_t\mathbf{u}+(\mathbf{u}\cdot\nabla)\mathbf{u}=-\frac{1}{\rho}\nabla p+\nu\nabla^2\mathbf{u}+\mathbf{f},\qquad \nabla\cdot\mathbf{u}=0,
\label{eq:ns}
\end{equation}
where $\mathbf{u}$ is velocity, $p$ is pressure, $\rho$ is density, $\nu$ is kinematic viscosity, and $\mathbf{f}$ is forcing. Subject to compatible changes of forcing, boundaries, and coordinates, Eq. \eqref{eq:ns} has translation, rotation, and Galilean symmetries. For a constant velocity $\mathbf{U}$, the Galilean change of frame is
\begin{equation}
\mathbf{x}'=\mathbf{x}-\mathbf{U}t,\qquad t'=t,\qquad \mathbf{u}'(\mathbf{x}',t')=\mathbf{u}(\mathbf{x},t)-\mathbf{U}.
\label{eq:galilean}
\end{equation}
These statements constrain physically admissible observables and numerical schemes. They do not imply that an increment PDF, a wall-normal profile, or a fitted exponent must coincide between a channel, a wake, a rotating flow, and forced homogeneous-isotropic turbulence (HIT).

This distinction matters for data analysis. An estimator may break an equation symmetry through its filter, coordinate convention, finite window, or sampling geometry. A claim of invariance must therefore name the object that is expected to remain unchanged: the equation, an ensemble law, a scale-normalized statistic, an exponent interval, a flux relation, or a learned representation. CURA-Turb records this as the first field of an invariance dossier.

\subsection{Exact conditional relations}
For a longitudinal velocity increment
\begin{equation}
\delta u_L(\mathbf{x},\mathbf{r})=\left[\mathbf{u}(\mathbf{x}+\mathbf{r})-\mathbf{u}(\mathbf{x})\right]\cdot\frac{\mathbf{r}}{r},
\label{eq:increment}
\end{equation}
the order-$p$ longitudinal structure function is
\begin{equation}
S_p(r)=\left\langle\left|\delta u_L(\mathbf{x},\mathbf{r})\right|^p\right\rangle.
\label{eq:structure}
\end{equation}
In a three-dimensional homogeneous, isotropic, stationary inertial-range setting, the celebrated third-order relation is
\begin{equation}
\left\langle\left[\delta u_L(r)\right]^3\right\rangle=-\frac{4}{5}\epsilon r,
\label{eq:fourfifths}
\end{equation}
where $\epsilon$ is the mean dissipation rate. Equation \eqref{eq:fourfifths} is not merely a fitted exponent: it is an exact conditional balance relation with assumptions that must travel with every use. Alternative exact-relation formulations can offer complementary cascade diagnostics \cite{andres_banerjee_2019}, but they do not remove the need to establish the physical and statistical conditions of the chosen relation.

\subsection{Similarity, universality, and intermittency}
The simplest K41 similarity form writes an inertial-range expectation as
\begin{equation}
S_p(r)\sim C_p(\epsilon r)^{p/3},
\label{eq:k41}
\end{equation}
where $C_p$ is an order-dependent coefficient. The relation is a useful baseline, not a guarantee that all orders or all flows have exponent $p/3$. A more general empirical representation is
\begin{equation}
S_p(r)\sim r^{\zeta_p},
\label{eq:zetap}
\end{equation}
with anomalous scaling indicated by $\zeta_p\neq p/3$ under the declared convention. Intermittency makes high-order moments particularly sensitive to rare events, finite sample support, and scaling-window selection. It is therefore possible for a low-order statistic to show a stable conditional regularity while a high-order statistic is intermittency modified.

The phrase ``ultimate invariance'' should not erase this hierarchy. A robust result might be an exact law under HIT assumptions, an exponent interval stable across a restricted family of HIT data sets, a distributional pattern that survives a specified filter, or a regime-dependent departure that identifies the physical influence of walls, shear, rotation, stratification, or decay. Each is scientifically valuable when its claim type and validity domain are explicit.

\section{Evidence Survey Across Regimes}
\subsection{Small scales and the conditional universality program}
Sreenivasan and Antonia review small-scale scaling, refined similarity, derivative statistics, intermittency models, and kinematics in a setting strongly connected to homogeneous turbulence \cite{sreenivasan_antonia_1997}. Frisch's synthesis likewise provides a cascade and multifractal language for why simple dimensional scaling can fail at higher orders \cite{frisch_1995}. These sources support a disciplined statement: the evidence motivates statistical regularity at small scales under appropriate separation and symmetry conditions, but it does not create a theorem that large-scale geometry or forcing becomes irrelevant in every observable at finite Reynolds number.

Onsager-inspired theory further emphasizes that cascade and dissipation questions have regularity and limiting-process content \cite{eyink_sreenivasan_2006}. This is relevant to invariance because a numerical or experimental statistic may be apparently scale invariant while still being sensitive to unresolved scales, viscosity, filtering, or nonstationary energy transfer. A universal-looking slope without an uncertainty and resolution audit is therefore weaker evidence than a typed exact-law check with a documented validity interval.

\subsection{Why flow families must not be pooled indiscriminately}
HIT provides a controlled reference because periodic forcing, stationarity, and isotropy can be explicitly targeted. Decaying turbulence removes the stationarity assumption and can shift the appropriate normalization. Wall-bounded turbulence introduces a distinguished direction and mean shear. Rotation and stratification introduce additional timescales and directional effects; an OpenAlex-discovered study by Sen \emph{et al.} motivates an anisotropy/nonuniversality control in rotating turbulence \cite{sen_2012}. Engineering flows add geometry, production, inhomogeneity, and measurement-volume effects. These cases may share a cascade vocabulary while differing in the meaning and observability of an ``inertial range.''

Time irreversibility is another example of a statistic that must be classified rather than romanticized. Vela-Martin and Jimenez connect inertial-range cascades to statistical irreversibility in a specified isotropic-turbulence setting \cite{vela_martin_jimenez_2021}. The existence of such a diagnostic does not make its value an invariant across every regime; it motivates a separate dossier field and a comparison only between cards with compatible temporal and sampling assumptions.

\subsection{Survey conclusion and gaps}
The Survey Agent accepts four gaps. First, exact symmetries, exact conditional relations, conditional similarity, and empirical cross-flow regularities are routinely conflated. Second, HIT evidence is too easily transferred to wall, shear, rotation, stratification, decaying, and measured flows. Third, finite scale windows, estimator choices, resolution, and sampling can alter apparent exponent or PDF collapse. Fourth, pointwise or low-order agreement does not establish high-order, distributional, or cross-flow invariance. These gaps motivate a protocol that records what a claim is before asking whether it is true.

\section{Research Gaps and Idea Selection}
The Idea Agent considered three directions. A single exponent catalogue is attractive because it permits a compact ranking, but it cannot identify whether a common slope is caused by a shared mechanism, a short fitting interval, a filter, or an inappropriate pooled ensemble. An exact-law-only benchmark is scientifically strong, particularly for Eq. \eqref{eq:fourfifths}, but it excludes important higher-order, anisotropic, and engineering-relevant questions. The selected direction is \emph{Conditional Universality and Regime Audit for Turbulence} (CURA-Turb).

CURA-Turb is a typed claim-and-evidence protocol. Each proposed invariant is represented by
\begin{equation}
\mathcal{D}=\left(\mathcal{T},Q,\mathcal{R},\mathcal{W},\mathcal{E},\mathcal{U},\mathcal{P}\right),
\label{eq:dossier}
\end{equation}
where $\mathcal{T}$ is claim type, $Q$ is the statistic, $\mathcal{R}$ is a regime card, $\mathcal{W}$ is a scale-window specification, $\mathcal{E}$ is an estimator, $\mathcal{U}$ is the uncertainty route, and $\mathcal{P}$ is the provenance/evidence tier. The claim types are \texttt{EQUATION\_SYMMETRY}, \texttt{EXACT\_CONDITIONAL\_LAW}, \texttt{CONDITIONAL\_SIMILARITY}, and \texttt{EMPIRICAL\_CROSS\_FLOW\_REGULARITY}. No type can be upgraded merely because a plot appears smooth.

Figure \ref{fig:cura} shows the workflow. The editable source is retained as \texttt{figures/cura\_turb\_flow.drawio}. The flow starts with a claim dossier, then tests assumptions, regime/estimator stability, and an evidence ledger. Its final label remains conditional and proposal-only.

\begin{figure*}[!t]
\centering
\fbox{\begin{minipage}{0.95\textwidth}
\centering\footnotesize
\textbf{CURA-Turb invariance audit}\\[5pt]
\begin{tabular}{c@{\quad$\longrightarrow$\quad}c@{\quad$\longrightarrow$\quad}c@{\quad$\longrightarrow$\quad}c}
\parbox{0.16\textwidth}{\centering\textbf{Claim dossier}\\symmetry\\exact law\\similarity\\empirical regularity} &
\parbox{0.16\textwidth}{\centering\textbf{Assumption gate}\\stationarity\\homogeneity\\isotropy\\scale interval\\representation} &
\parbox{0.18\textwidth}{\centering\textbf{Regime and estimator stress}\\HIT, decay, wall\\rotation, shear\\measurement\\window stability} &
\parbox{0.17\textwidth}{\centering\textbf{Evidence ledger}\\law residual\\exponent interval\\PDF distance\\anisotropy\\uncertainty}\\[14pt]
\multicolumn{4}{c}{\textbf{Status:} exact-law eligible, conditional universality, intermittency modified, regime dependent, inconclusive, or invalid.}
\end{tabular}
\end{minipage}}
\caption{CURA-Turb treats universality as a claim with assumptions and evidence requirements. The figure is a research-design artifact, not an observed flow classification.}
\label{fig:cura}
\end{figure*}

\section{ExperimentDesign: CURA-Turb Computational Protocol}
\subsection{Regime cards and compatibility}
The ExperimentDesign Agent selects the computational-digital template. Its unit of analysis is one \emph{claim dossier $\times$ regime card $\times$ scale-window card $\times$ estimator card $\times$ evidence card}. Table \ref{tab:cards} specifies seven cards. `F0` holds analytic equation-symmetry and balance references. `F1` is forced periodic HIT. `F2` is decaying nominally isotropic turbulence. `F3` represents wall-bounded flows. `F4` represents rotating or stratified turbulence. `F5` represents sheared, inhomogeneous, or engineering flows. `F6` records experimental or observational representation. The cards are not a taxonomy for its own sake: they prevent a statement from losing the assumptions that give it meaning.

\begin{table*}[!t]
\caption{CURA-Turb regime cards and validation obligations. All fields are planned inputs or derived analysis records; this proposal supplies no observed data.}
\label{tab:cards}
\centering
\small
\begin{tabularx}{\textwidth}{lYYY}
\toprule
Card & Purpose & Required fields & Failure response \\
\midrule
F0 analytic reference & Separate equation symmetry from measured statistical claim & equation, transformation, forcing/boundary assumptions & \texttt{MODEL\_INVALID} if assumptions conflict \\
F1 forced periodic HIT & Eligible reference for conditional exact-law tests & stationarity, homogeneity/isotropy diagnostics, dissipation route, scale separation & no transfer without a bridge \\
F2 decaying turbulence & Track nonstationary normalization and time origin & decay protocol, time marker, scale interval & \texttt{REGIME\_DEPENDENT} or inconclusive \\
F3 wall-bounded flow & Preserve wall distance, mean shear, and inhomogeneity & wall coordinate, shear scaling, sampling geometry & prohibit isotropic pooling \\
F4 rotating/stratified flow & Preserve directional and extra-timescale effects & Rossby/Froude context, directional statistics, anisotropy marker & \texttt{REGIME\_DEPENDENT} \\
F5 engineering/sheared flow & Audit geometry, production, and transfer bridge & forcing, geometry, production, measurement volume & withhold HIT universality claim \\
F6 measured representation & Harmonize data and estimator before comparison & instrument/filter, resolution/noise, cadence, reconstruction & \texttt{REPRESENTATION\_INVALID} \\
\bottomrule
\end{tabularx}
\end{table*}

Compatible comparisons are not necessarily identical flows. For example, multiple `F1` cases may be compared after their scale-separation, increment convention, and dissipation estimates are made compatible. An `F1` result may be placed beside an `F3` result only as a declared transfer test with wall-distance and shear markers, not as an automatic pooled estimate. A physical statistic can recur across cards while the status of its invariance claim changes.

\subsection{Metrics, estimator audit, and uncertainty ledger}
For an exact-law dossier, the planned residual can be written conceptually as
\begin{equation}
R_{4/5}(r)=\frac{\left\langle[\delta u_L(r)]^3\right\rangle+(4/5)\epsilon r}{\epsilon r},
\label{eq:residual}
\end{equation}
to be evaluated only in an eligible `F1`-like setting. The residual is not a universal quality score. It requires an explicitly estimated scale interval, dissipation route, uncertainty interval, and statistical assumptions. For conditional similarity, CURA-Turb records exponent intervals rather than a single slope. For distributional comparisons, it records a specified normalized PDF distance and the normalization, filter, and sampling choices that define it.

The protocol uses a conservative claim ledger,
\begin{align}
\epsilon_Q^{\mathrm{claim}} &\lesssim \epsilon_{\mathrm{model}}+\epsilon_{\mathrm{repr}}+\epsilon_{\mathrm{sampling}}+\epsilon_{\mathrm{window}} \nonumber\\
&+\epsilon_{\mathrm{estimator}}+\epsilon_{\mathrm{regime}}+\epsilon_{\mathrm{transfer}},
\label{eq:errorledger}
\end{align}
where the terms represent model/discretization limitations, data representation, finite sampling, scale-window choice, estimator form, regime mismatch, and cross-card transfer. Equation \eqref{eq:errorledger} is an evidence ledger, not an independent-error identity. If a term is unknown, the output must retain it as an uncertainty or return an inconclusive status.

The planned outcome vector is
\begin{align}
\mathbf{y}=\bigl(&R_{4/5},[\zeta_p]_{\mathrm{interval}},d_{\mathrm{PDF}},A_{\mathrm{aniso}},I_{\mathrm{irr}}, \nonumber\\
&\epsilon_Q^{\mathrm{claim}},\sigma_{\mathrm{status}}\bigr),
\label{eq:outcomevector}
\end{align}
where $[\zeta_p]_{\mathrm{interval}}$ is a structure-function exponent interval, $d_{\mathrm{PDF}}$ is a declared distribution-distance measure, $A_{\mathrm{aniso}}$ is an anisotropy/inhomogeneity marker, $I_{\mathrm{irr}}$ is a time-irreversibility marker, and $\sigma_{\mathrm{status}}$ records sensitivity of the final label to allowed estimator choices. No single component is permitted to dominate the conclusion by default.

\subsection{Validation ladder and falsifiers}
The first gate verifies units, coordinate frames, increment definitions, filtering, sampling, and data provenance. The second gate records the assumptions needed by the claim. A four-fifths test without a declared stationarity, homogeneity, isotropy, and inertial-range route cannot receive even the label \texttt{EXACT\_RELATION\_ELIGIBLE}. The third gate applies a predeclared estimator audit: scale windows, fitting formulations, sampling blocks, and compatible normalizations are varied only within documented bounds. The fourth gate performs a regime stress test using cards that can answer the intended transfer question. The fifth gate requires an explicit bridge before an idealized result is written as an engineering statement.

The primary idea fails if the classification is not more stable and interpretable than an exponent-only comparison. It also fails if the same dossier can move from exact law to universality claim without new evidence, or if an apparent cross-flow agreement is explained entirely by changed filtering or a shared finite fitting interval. These falsifiers make the audit scientifically useful even when it returns a negative result.

\section{Expected Outcome Branches and Conditional Conclusions}
Table \ref{tab:branches} defines authorized output language. The purpose is not to make results cautious by default. It permits a direct conclusion when a claim has earned one, while reserving the word ``universal'' for a conditionally stable comparison across eligible cards. A favorable label never represents a new simulation or measurement in this proposal; it describes what a future reviewed execution would be allowed to conclude.

\begin{table}[!t]
\caption{CURA-Turb conditional-status framework. Each row is a predeclared interpretation, not a reported result.}
\label{tab:branches}
\centering
\footnotesize
\begin{tabularx}{\columnwidth}{lY}
\toprule
Label & Authorized conditional interpretation \\
\midrule
\texttt{EXACT\_RELATION\_ELIGIBLE} & The stated physical and statistical conditions support testing the named exact relation; numerical agreement remains to be evaluated. \\
\texttt{CONDITIONALLY\_UNIVERSAL} & The dossier statistic is stable across eligible matched cards and estimator intervals; the validity domain remains explicit. \\
\texttt{INTERMITTENCY\_MODIFIED} & Higher-order or distributional evidence departs from the simplest similarity model within the declared comparison. \\
\texttt{REGIME\_DEPENDENT} & Direction, wall, shear, rotation, stratification, forcing, or nonstationarity changes the claim. \\
\texttt{REPRESENTATION\_INVALID} & Units, filtering, sampling, reconstruction, or provenance fails the dossier contract. \\
\texttt{INCONCLUSIVE} & Scale separation, sample support, estimator stability, or transfer evidence is insufficient. \\
\texttt{MODEL\_INVALID} & Equation assumptions, regime assignment, or observation/projection status is incompatible. \\
\bottomrule
\end{tabularx}
\end{table}

The branches directly answer the topic. The strongest candidates for invariance are not a single empirical exponent but a layered set of statements: exact equation symmetries, exact conditional balance laws, and conditional low-scale statistical regularities within compatible regimes. Intermittency and regime dependence are not failures of turbulence science; they are the information that limits and sharpens the validity domain. CURA-Turb makes those distinctions visible enough to be tested rather than buried in prose.

\section{Pre-Registered Decision Audit}
A robust universality assessment must declare its decision logic before a convenient scaling window, filtered signal, or subset of flow records is selected. CURA-Turb therefore registers the statistic, claim type, eligible cards, normalization, scale interval, estimator family, uncertainty method, and transfer target before analysis. This prevents a familiar failure mode: a visually persuasive log-log line is fitted after looking at the data and then treated as evidence that a theory holds independently of its assumptions.

Table \ref{tab:audit} lists the required audit. It distinguishes a missing quantity from a negative physical result. A dossier with an unverified increment convention is not an example of nonuniversal turbulence; it is representation invalid. A high-order fit that changes across reasonable windows is not evidence for a universal exponent; it is inconclusive or intermittency modified. A wall-flow statistic that differs from an HIT reference is not a failed experiment; it can be a regime-dependent result that identifies where the physical mechanism matters.

\begin{table*}[!t]
\caption{Pre-registered CURA-Turb decision audit. All entries are future evidence obligations and do not assert an executed calculation.}
\label{tab:audit}
\centering
\small
\begin{tabularx}{\textwidth}{lYYY}
\toprule
Audit item & Required record & Risk prevented & Authorized response if incomplete \\
\midrule
Claim identity & Claim type, statistic, increment/normalization convention, physical dimension and units & An exact relation, a fit, and a symmetry are conflated & \texttt{MODEL\_INVALID} \\
Assumption identity & Stationarity, homogeneity, isotropy, incompressibility, forcing/boundary route, and scale eligibility & Four-fifths or K41 language used outside its scope & Restrict claim or mark ineligible \\
Representation identity & Instrument/filter, reconstruction, DNS discretization, cadence, and coordinate frame & Measurement or numerical representation is mistaken for flow physics & \texttt{REPRESENTATION\_INVALID} \\
Estimator stability & Predeclared windows, fit model, uncertainty interval, resampling/block rule, and alternative estimator & One convenient slope becomes a universal exponent & \texttt{INCONCLUSIVE} \\
Regime stress & F1--F6 compatibility, direction/wall/rotation/shear markers, and matching rule & HIT evidence is silently pooled with an incompatible flow & \texttt{REGIME\_DEPENDENT} \\
Transfer bridge & Target application, missing physics, domain expert review, and observation/projection flag & Idealized result is retold as an engineering guarantee & Withhold engineering transfer \\
\bottomrule
\end{tabularx}
\end{table*}

The audit makes negative findings productive. A conditional HIT regularity can validate a baseline for simulation verification. A regime-dependent label can identify the role of walls, rotation, or production in an engineering model. An intermittency-modified label can specify which orders or tails need a non-K41 closure. An inconclusive label can reveal whether additional scale separation, better temporal/spatial resolution, or a better uncertainty model is the real bottleneck. The goal is not to eliminate variation, but to say which variation is physical, which is representational, and which remains unresolved.

\section{Engineering and Computational Implications}
\subsection{A validity-domain language for models}
Engineering turbulence models are often evaluated on flows that differ from their calibration conditions. CURA-Turb does not replace a closure model, a large-eddy simulation, a direct numerical simulation, or a measurement campaign. It provides a compact language for documenting whether a statistical regularity used in such a model is tied to an exact balance, a calibrated conditional similarity, or a cross-regime extrapolation. The distinction can guide model validation: an exact-law residual is a physics-consistency check; an exponent interval is a conditional scaling diagnostic; a regime marker is a transfer-risk signal.

This distinction is useful for machine learning as well. A learned surrogate can show impressive aggregate accuracy while violating Galilean consistency, depending on a filter, or exploiting a regime imbalance in its training set. A CURA-Turb dossier gives a learned model testable invariance obligations: declare frame transformation, scale representation, regime coverage, and an uncertainty/transfer route. It does not presume that an invariant neural representation exists; it makes that hypothesis falsifiable.

\subsection{A bridge to complex nonlinear systems}
Turbulence is an archetype of a complex nonlinear system because it combines nonlinear coupling, a broad range of scales, intermittency, and sensitivity to constraints. The proposal's claim taxonomy can transfer cautiously to other systems: distinguish a governing-equation symmetry from an emergent ensemble regularity; state the coarse-graining and sampling rules; and test whether a result is stable across compatible regimes. The transfer is conceptual, not an assertion that all complex systems share Navier--Stokes cascade laws.

\section{Risks, Limitations, and Human Review Requirements}
The report deliberately does not resolve the mathematical existence and regularity problem for three-dimensional Navier--Stokes equations. It also does not establish an infinite-Reynolds-number limit, compute a new exponent, or settle a universality debate. The K41 form in Eq. \eqref{eq:k41}, the refined similarity context, and intermittency phenomenology all require careful convention and regime handling. The exact four-fifths relation in Eq. \eqref{eq:fourfifths} is a stringent reference only where its stated assumptions are justified.

Finite data and finite computation are major risks. A DNS may lack a clean scale interval; an experiment may have spatial/temporal resolution and noise constraints; an engineering flow may lack homogeneity or isotropy; a rotating or stratified flow may require direction-resolved statistics. An AI system can make these risks worse by producing a clean narrative from incompatible results. CURA-Turb counteracts that failure mode through required cards and refusal statuses, but it cannot supply missing evidence.

The strongest directly checked evidence here is the publisher-verified Sreenivasan--Antonia review \cite{sreenivasan_antonia_1997}. Other background and discovery sources remain subject to publisher/DOI-page review before external publication. Any implementation needs turbulence-theory, numerical-analysis, statistics/uncertainty-quantification, and experimental-metrology review; an engineering application also needs a domain flow expert. The execution policy is \texttt{DESIGN\_ONLY}: no simulator run, laboratory acquisition, facility operation, or observed data are represented as outputs of this work.

\section{Conclusion}
The ultimate statistical invariance of turbulence is best understood as a hierarchy rather than a single number. Governing equations possess exact structural symmetries. Certain inertial-range relations, such as the four-fifths law, are exact under well-defined assumptions. Similarity and small-scale regularity can be conditionally robust in compatible regimes. High-order intermittency, anisotropy, boundaries, shear, rotation, stratification, nonstationarity, finite resolution, and estimator choices can alter the apparent statistic or limit its transfer.

CURA-Turb converts that answer into a research design. It prevents an exact law from being mistaken for an empirical collapse, prevents an exponent fit from becoming a universal theorem, and prevents an idealized HIT result from silently becoming an engineering guarantee. Its evidence ledger, regime cards, estimator stress tests, and conditional statuses allow the research program to make strong claims where the evidence supports them and useful negative claims where it does not. The next step is a specialist-reviewed implementation on carefully selected turbulence data sets or simulations with predeclared dossiers and transparent uncertainty reporting.

\appendices
\section{Evidence Coverage and Claim Boundary}
The Survey registry contains eight bounded sources. The Sreenivasan--Antonia review was checked directly on the Annual Reviews publisher page \cite{sreenivasan_antonia_1997}. Kolmogorov, Frisch, and Eyink--Sreenivasan sources provide background for exact relations, similarity, intermittency, and cascade regularity \cite{kolmogorov_1941,kolmogorov_1962,frisch_1995,eyink_sreenivasan_2006}. Vela-Martin--Jimenez, Sen \emph{et al.}, and Andres--Banerjee are discovery-stage sources used only for the scoped diagnostic motivations stated in the Survey \cite{vela_martin_jimenez_2021,sen_2012,andres_banerjee_2019}. No source is used to claim a new DNS run, new laboratory measurement, universal exponent, or engineering-flow prediction.

\section{Symbols and Operational Definitions}
\begin{tabularx}{0.96\columnwidth}{lY}
\toprule
Symbol / term & Definition in this proposal \\
\midrule
$\mathbf{u}$, $p$, $\rho$, $\nu$, $\mathbf{f}$ & Velocity, pressure, density, viscosity, and forcing in the scoped incompressible equation. \\
$\delta u_L$, $S_p(r)$ & Longitudinal increment and order-$p$ structure function in Eqs. \eqref{eq:increment}--\eqref{eq:structure}. \\
$\epsilon$, $R_{4/5}$ & Mean dissipation convention and declared four-fifths residual diagnostic. \\
$\zeta_p$ & Fitted or estimated scaling-exponent interval, not automatically a universal constant. \\
$d_{\mathrm{PDF}}$ & Predeclared normalized probability-distribution distance. \\
$A_{\mathrm{aniso}}$, $I_{\mathrm{irr}}$ & Anisotropy/inhomogeneity and time-irreversibility markers. \\
$\epsilon_Q^{\mathrm{claim}}$ & Conservative evidence/error ledger for a claim about statistic $Q$. \\
CURA-Turb & Conditional Universality and Regime Audit for Turbulence. \\
\bottomrule
\end{tabularx}

\section{Minimum Reproducible Claim-Dossier Schema}
The proposed output is not a generic ``universality score.'' It is a compact, reviewable dossier that can be handed from a turbulence analyst to a numerical, experimental, or engineering reviewer without silently changing the meaning of a claim. Every dossier must first name a \emph{claim object}: an equation symmetry, an exact conditional law, a similarity statement, or an empirical cross-flow regularity. It then stores the physical statistic, the increment convention, the normalization, the eligible regime cards, the scale interval, and the estimator family. A statistic without these records can still be scientifically interesting, but it cannot receive an invariance label under this proposal.

The record also separates three kinds of incompleteness. \emph{Physical incompleteness} occurs when a condition such as isotropy, stationarity, or a needed forcing route is not established. \emph{Representation incompleteness} occurs when a filter, reconstruction, coordinate system, sampling cadence, or numerical discretization is unknown or incompatible. \emph{Inferential incompleteness} occurs when the scale interval, sample support, fitting rule, or uncertainty method is not adequate to distinguish compatible conditional explanations. These cases must not be collapsed into a single low confidence value: they lead to different follow-up actions and different authorized status labels.

\begin{table}[!t]
\caption{Minimum reproducible CURA-Turb claim dossier. Fields are required future records, not reported measurements.}
\label{tab:schema}
\centering
\footnotesize
\begin{tabularx}{\columnwidth}{lY}
\toprule
Field & Required content and release rule \\
\midrule
Claim object & One of the four claim types; a plain-language sentence that can be falsified. \\
Statistic definition & Component, increment direction, order, normalization, units, and sign convention. \\
Eligibility contract & Regime card, symmetry/balance assumptions, scale-separation route, and excluded cases. \\
Representation record & Data or simulation provenance, coordinates, filtering, resolution, sampling, and reconstruction. \\
Estimator plan & Predeclared scale window, fit or residual form, alternatives, and resampling/block strategy. \\
Uncertainty route & Model, representation, sampling, window, estimator, regime, and transfer uncertainties. \\
Comparison map & Matched cards, matching criteria, deliberately unmatched cards, and transfer target. \\
Release decision & Conditional status, reasons, unresolved fields, human-review roles, and permitted wording. \\
\bottomrule
\end{tabularx}
\end{table}

The release decision is deliberately asymmetric. A complete eligible dossier may authorize the bounded statement that an exact relation is ready to be tested, that a conditional similarity remains stable over a declared matched comparison, or that an observed departure would be classified as intermittency modified or regime dependent. It never authorizes the stronger sentence that all turbulence has been shown to share a statistic. Conversely, an incomplete dossier must say \texttt{INCONCLUSIVE}, \texttt{REPRESENTATION\_INVALID}, or \texttt{MODEL\_INVALID} rather than report a negative universality result. This distinction preserves the scientific value of a carefully documented failure to transfer.

The four-agent handoff is explicit. The Survey Agent contributes the source key, evidence tier, claim boundary, and gap identifier. The Idea Agent contributes the selected comparison and falsifier: for example, whether a candidate relation should survive compatible F1 cards but be stress-tested against F3--F5. The ExperimentDesign Agent adds the estimator and uncertainty contract, while the Author Agent may write only the status and conclusion supported by those fields. The handoff prevents a later narrative from upgrading a discovery-stage reference into direct evidence or changing a conditional benchmark into an executed result.

For a future implementation, dossiers should be released in machine-readable and human-readable forms. A machine-readable record supports consistent audits across many velocity fields, while the human-readable record must state why cards are comparable, what physics has been excluded, and which specialist must review the status. At minimum, a turbulence theorist reviews the claim type and assumptions; a numerical or metrology expert reviews representation and uncertainty; and a domain-flow expert reviews transfer to an engineering application. Such review is a condition of use, not a substitute for the evidence ledger. The present paper provides this schema only; it creates no dossier from new flow data and performs no numerical or laboratory evaluation.

\begin{thebibliography}{99}
\bibitem{sreenivasan_antonia_1997} K. R. Sreenivasan and R. A. Antonia, ``The phenomenology of small-scale turbulence,'' \emph{Annu. Rev. Fluid Mech.}, vol. 29, pp. 435--472, 1997, doi: 10.1146/annurev.fluid.29.1.435.
\bibitem{kolmogorov_1941} A. N. Kolmogorov, ``The local structure of turbulence in incompressible viscous fluid for very large Reynolds numbers,'' \emph{Dokl. Akad. Nauk SSSR}, vol. 30, pp. 299--303, 1941; English translation, \emph{Proc. R. Soc. Lond. A}, vol. 434, pp. 9--13, 1991, doi: 10.1098/rspa.1991.0075.
\bibitem{kolmogorov_1962} A. N. Kolmogorov, ``A refinement of previous hypotheses concerning the local structure of turbulence in a viscous incompressible fluid at high Reynolds number,'' \emph{J. Fluid Mech.}, vol. 13, pp. 82--85, 1962, doi: 10.1017/S0022112062000518.
\bibitem{frisch_1995} U. Frisch, \emph{Turbulence: The Legacy of A. N. Kolmogorov}. Cambridge, U.K.: Cambridge Univ. Press, 1995.
\bibitem{eyink_sreenivasan_2006} G. L. Eyink and K. R. Sreenivasan, ``Onsager and the theory of hydrodynamic turbulence,'' \emph{Rev. Mod. Phys.}, vol. 78, pp. 87--135, 2006, doi: 10.1103/RevModPhys.78.87.
\bibitem{vela_martin_jimenez_2021} A. Vela-Martin and J. Jimenez, ``Entropy, irreversibility and cascades in the inertial range of isotropic turbulence,'' \emph{J. Fluid Mech.}, vol. 915, 2021, doi: 10.1017/jfm.2021.105.
\bibitem{sen_2012} A. Sen, P. D. Mininni, D. Rosenberg, and A. Pouquet, ``Anisotropy and nonuniversality in scaling laws of the large-scale energy spectrum in rotating turbulence,'' \emph{Phys. Rev. E}, vol. 86, 036319, 2012, doi: 10.1103/PhysRevE.86.036319.
\bibitem{andres_banerjee_2019} N. Andres and S. Banerjee, ``Statistics of incompressible hydrodynamic turbulence: An alternative approach,'' \emph{Phys. Rev. Fluids}, vol. 4, 024603, 2019, doi: 10.1103/PhysRevFluids.4.024603.
\end{thebibliography}
\end{document}
